\section{Review of differentiable manifolds}

\begin{definition}[n-dimensional differentiable manifold]
An \emph{n-dimensional differentiable manifold} is a topological space $M$ 
with an open cover $\{U_{\alpha}\}_{\alpha}$ and maps
\[\{\phi_{\alpha} : U_{\alpha} \to \phi_{\alpha}(U_{\alpha})\subset\mathbb{R}^n\}_{\alpha \in \Lambda}\]
satisfying
\begin{enumerate}
    \item $\phi_{\alpha}(U_{\alpha})\subset\mathbb{R}^n$ is open and $\phi_{\alpha}:U_{\alpha} \to \phi_{\alpha}(U_{\alpha})$ is a homeomorphism.
    \item If $U_{\alpha}\cap U_{\beta}\neq \emptyset$, then the transition map $\phi_{\beta} \circ \phi_{\alpha}^{-1}$ is a diffeomorphism (between open subsets of $\mathbb{R}^n$).
\end{enumerate}
$\{(U_{\alpha}, \phi_{\alpha}) : \alpha \in \Lambda\}$ is called a \emph{coordinate neighborhood system} or \emph{atlas}.
\end{definition}

Some terminology
\begin{itemize}
\item $n = \operatorname{dim}M$ \emph{dimension} of $M$.
\item $\phi_{\alpha}$ \emph{coordinate map} or \emph{chart}.
\item $\phi_{\alpha} = (x^1, \dots, x^n)$ \emph{local coordinates} $x^i : U_{\alpha} \to \mathbb{R}$.
\end{itemize}

\begin{definition}[Compatible charts]
Assume $M$ is an n-differentiable manifold with atlas $\{(U_{\alpha}, \phi_{\alpha}) : \alpha \in \Lambda\}$. A chart $(U, \phi)$ is \emph{compatible} with the atlas if
$\{U_{\alpha}, \phi_{\alpha})\}_{\alpha \in \Lambda} \cup (U,\phi)$ is again an atlas.
The set of all such charts is called the \emph{maximal coordinate neighborhood system} (or \textit{atlas}) of $M$.
\end{definition}

Henceforth, we assume the atlases of $M$ are always maximal.

\subsection{Differential of maps and tangent spaces}

\begin{definition}[$C^{\infty}$]
Assume $M,N$ are differentiable manifolds. A map continuous $f:M\to N$ is \emph{$C^{\infty}$} at $p \in M$
if there exists charts $(U, \phi)$ around $p$ and $(V, \psi)$ around $f(p)$ such that $\psi \circ f \circ \phi^{-1}: \phi(U) \to \psi(V)$ is $C^{\infty}$.
\end{definition}

\begin{remark}
\begin{itemize}
\item $M = I \subset \mathbb{R}, c:I\to N$ is a $C^{\infty}$ curve.
\item $N = \mathbb{R}, f:M\to \mathbb{R}$ is a $C^{\infty}$ function.
\item $C^{\infty}(M):= \{C^{\infty}\text{ function on } M\}$.
\end{itemize}
\end{remark}

\begin{definition}[Tangent vector and tangent space]
A map $v: C^{\infty}(M) \to \mathbb{R}$ is a \emph{tangent vector} at $p \in M$ if for all $a,b\in\mathbb{R}$ and all $f,g \in C^{\infty}(M)$, we have
\begin{enumerate}
\item $v(af + bg) = av(f) + bv(g)$ (linearity).
\item $v(fg) = f(p)v(g) + g(p)v(f)$ (Leibniz rule).
\end{enumerate}
The set $T_pM : = \{v : v \text{ is a tangent vector at } p\}$ is the \emph{tangent space} to $M$ at $p$.
\end{definition}

\begin{definition}[Velocity vector]
Let $c:I \to M$ be a $C^{\infty}$ curve with $c(0) = p \in M$. The \emph{velocity vector} of $c$ at $p$ is the map $c'(0) \in T_pM$ defined by
\[c'(0)(f) = \frac{d}{dt}\bigg|_{t=0} (f \circ c)(t), \quad f \in C^{\infty}(M).\]
\end{definition}

\begin{definition}[Coordinate vector fields]
Let $(U,\phi)$ be a chart around $p$ with $\phi(p)=0$, $\phi=(x^1,\dots,x^n)$.
Define curves
\[
c_i(t) := \phi^{-1}(0,\dots,t,\dots,0).
\]
Then
\[
\left.\frac{\partial}{\partial x^i}\right|_p := c_i'(0).
\]
\end{definition}

\begin{theorem}
For $p \in M$ an n-differentiable manifold, the tangent space $T_pM$ is an n-dimensional vector space with basis $\{\frac{\partial}{\partial x^1}\big|_p, \dots, \frac{\partial}{\partial x^n}\big|_p\}$.
I.e.
\[T_pM = \operatorname{span}\left\{\left.\frac{\partial}{\partial x^1}\right|_p, \dots, \left.\frac{\partial}{\partial x^n}\right|_p\right\}.\]
\end{theorem}

\begin{definition}[Differential]
Let $f:M\to N$ be a $C^{\infty}$ map between differentiable manifolds. The \emph{differential} of $f$ at $p \in M$ is the map
$df_p: T_pM \to T_{f(p)}N$ defined by
\[df_p(v)(g) := (f\circ c)'(0), \quad c \in C^{\infty}, \quad c(0) = p \text{ and } c'(0) = v. \]
\end{definition}

\begin{remark}
$c$ in the above definition is not unique, but $df_p(v)$ is well-defined.
\end{remark}

\begin{definition}[Diffeomorphism]
Let $M,N$ be differentiable manifolds. $f:M\to N$ is a \emph{diffeomorphism} if
\begin{itemize}
\item $f$ is a bijection.
\item $f$ and $f^{-1}$ are $C^{\infty}$.
\end{itemize}
In this case, $M$ and $N$ are \emph{diffeomorphic}.
\end{definition}

\subsection{Diffeomorphism groups and quotient manifolds}

Let $M$ be an n-differentiable manifold.

\begin{definition}[$C^{\infty}$ transformation group]
A set of diffeomorphisms $\Gamma = {\gamma_i}_{i\in I}$ on $M$ is a \emph{$C^{\infty}$ transformation group} if
\begin{itemize}
    \item For all $\gamma_i, \gamma_j \in \Gamma$, we have $\gamma_i \circ \gamma_j \in \Gamma$.
    \item For all $\gamma_i \in \Gamma$, we have $\gamma_i^{-1} \in \Gamma$.
    \item The identity map $\operatorname{Id}_M$ is in $\Gamma$.
\end{itemize}
\end{definition}

\begin{definition}[Proper discontinuous transformation group]
A $C^{\infty}$ transformation group $\Gamma$ on $M$ is \emph{properly discontinuous} if for each $p \in M$, there exists a chart $(U, \phi)$ around $p$ such that 
\[\gamma_i \in \Gamma \setminus \{\operatorname{Id}_M\} \implies \gamma_i(U) \cap U = \emptyset \text{ for all } i \in I.\]
\label{def:properly_disc_transf_group}
\end{definition}

Let $\Gamma$ be a properly discontinuous transformation group on $M$.
For each $p \in M$, let $(U_p, \phi_p)$ be a chart around $p$ as in definition \ref{def:properly_disc_transf_group}.

\begin{lemma}
For all $p\in M$ and all $\gamma_i, \gamma_j \in \Gamma$ ($\gamma_i \neq \gamma_j$), we have $\gamma_i(U_p) \cap \gamma_j(U_p) = \emptyset$ or $\gamma_i(U_p) = \gamma_j(U_p)$.
\end{lemma}

\begin{proof}
    $\gamma_i \neq \gamma_j \implies \gamma_i^{-1}\circ\gamma_j\neq\operatorname{Id}_M$. Then, by defintition \ref{def:properly_disc_transf_group},
    $U_p \cap \gamma_i^{-1}\circ\gamma_j(U_p) = \emptyset \implies \gamma_i(U_p) \cap \gamma_j(U_p) = \emptyset$.
\end{proof}

\begin{definition}[Orbit equivalence relation]
Let $\Gamma$ be a group acting on $M$. Define a relation on $M$ by
\[
p \sim q \iff \exists \gamma \in \Gamma \text{ such that } q = \gamma(p).
\]
\end{definition}

\begin{proposition}
The relation $\sim$ is an equivalence relation.
\end{proposition}

\begin{definition}[Orbit]
The equivalence class of $p$,
\[
[p] = \{ \gamma(p) \mid \gamma \in \Gamma \} \in M/\sim,
\]
is  the \emph{orbit} of $p$.
\end{definition}

\begin{definition}[Quotient manifold]
The quotient set $M/\Gamma := M/\sim$ is called the \emph{quotient manifold} of $M$ by $\Gamma$.
\begin{itemize}
    \item $M/\Gamma$ is equipped with the quotient topology via the projection map $\pi: M \to M/\Gamma$.
    \item $D \subset M$ is a \emph{fundamental domain} if $\pi|_D:D \to M/\Gamma$ is surjective.
\end{itemize}
\end{definition}

To show that $M/\Gamma$ is a differentiable manifold, we construct charts.
For $p \in M$, let $\hat{U}_p := \pi(U_p) \subset M/\Gamma$.

\begin{lemma}
If $q, r \in U_p$ satisfy $q \sim r$ then $q = r$. In particular, $\pi|_{U_p}: U_p \to \hat{U}_p$ is a
bijection.
\end{lemma}

\begin{proof}
$q \sim r \implies \exists \gamma_i \in \Gamma : r = \gamma_i(q) \in \gamma_i(U_p)$.
Thus $r \in \gamma_i(U_p)$ for some $\gamma_i \in \Gamma$.
Since $\Gamma$ is properly discontinuos, $\gamma_i = \operatorname{Id}_M$.
Hence, $r = \gamma_i(q) = q$. The second claim follows from $\pi(q) = \pi(r) \iff q \sim r$.
\end{proof}

\begin{problem}
Define $\hat{\phi}:\hat{U}_p \to \hat{\phi}_p(\hat{U}_p) = \phi_p(U_p) \subset \mathbb{R}^n$ by
\[\hat{\phi}_p := \phi_p \circ (\pi|_{U_p})^{-1}.\]
Show that $\hat{\phi}_p$ is a homeomorphism.
(Hint: In general, quotient maps are not necessarily open; use the proply discontinous property).
\end{problem}

\begin{theorem}
$\{(\hat{U}_p, \hat{\phi}_p) : p \in M\}$ is an atlas for $M/\Gamma$. Hence, $M/\Gamma$ is a differentiable manifold.
\end{theorem}

\subsection{Topology of differentiable manifolds}

Recall the convergence of a sequence $\{p_i\}_{i=1}^{\infty}$:
\[\lim_{i\to \infty} p_i = p \iff \forall \text{ open } O \ni p, \exists i_0 \in \mathbb{N} : i > i_0 \implies p_i \in O.\]

\begin{proposition}[Convergence in differentiable manifolds]
$\lim_{i\to\infty} p_i = p$ if and only if for every chart $(U, \phi)$ around $p$, we have $\lim_{i\to\infty} \phi(p_i) = \phi(p)$ (in $\mathbb{R}^n$).
That is, 
\[\lim_{i\to\infty} p_i = p \iff \forall (U, \phi) \exists i_0 \in \mathbb{N} : i > i_0 \implies \phi(p_i) \to \phi(p).\]
\end{proposition}

\begin{remark}
$M$ is Haussdorff and second countable and locally Euclidean. 
Thus, $f: M \to N$ is continuous at $p \in M$ if and only if $p_i \to p$ implies $f(p_i) \to f(p)$.
\end{remark}

\begin{lemma}
Let $(U,\phi)$ be a chart, $c:(a,b) \to M$ be a $C^{\infty}$ curve and $t_0 \in (a,b)$ that satisfy 
\begin{enumerate} 
    \item The map $\phi \circ c: (a,b) \to \mathbb{R}^n$ is $C^{\infty}$.
    \item $\exists \{t_i\}_{i=1}^{\infty} \subset (a,b) : t_i \to t_0$ and $t_i \geq t_0$ and $c(t_i) \notin U$.
\end{enumerate}
Then
\begin{enumerate}
    \item $c(t_0) \notin U$.
    \item $\forall \{s_i\}_{i=1}^{\infty} \subset (a,b) : s_i \to t_0$ and $s_i < t_0$, $(\phi \circ c)(s_i)$ does not converge in $\phi(U)$.
    \label{num}
\end{enumerate}
\end{lemma}

\begin{proof}
(1) Assume $c(t_0) \in U$. 
By continuity, $\exists \delta_0 > 0$ such that $c((t_0 - \delta_0, t_0 + \delta_0)) \subset U$.
For sufficiently large $i$, we have $t_i \in (t_0, t_0 + \delta_0) \implies c(t_i) \in U$, which contradicts \ref{num}.
(2) Assume $\exists \{s_i\} \to t_0$ ($s_i < t_0$) such that $\{\phi \circ c(s_i)\}$ converges in $\phi(U)$.
Then $(\phi \circ c)(s_i) \to x_0$ for some $x_0 \in \phi(U)$. By continuity of $\phi^{-1}$,
$c(s_i) = (\phi^{-1} \circ \phi)(c(s_i)) \to \phi^{-1}(x_0)$ as $i \to \infty$. By the continuity of $c$, we have $c(s_i) \to c(t_0)$ as $i \to \infty$. 
Since $M$ is Haussdorff, the limit is unique, so $\phi^{-1}(x_0) = c(t_0)$. Thus, $c(t_0) = \phi^{-1}(x_0) \in U$, which contradicts (1).
\end{proof}

\subsection{Tangent bundle}

Let $M$ be a differentiable manifold.

\begin{definition}[Tangent bundle]
The \emph{tangent bundle} of $M$ is the set
\[TM := \bigsqcup_{p \in M} T_pM = \{(p,v) : p \in M, v \in T_pM\}.\]
\end{definition}

We introduce a differentiable structure on $TM$ using the following framework:

\begin{problem}
Let $\{X_{\lambda}, \mathcal{O}_{\lambda}\}_{\lambda \in \Lambda}$ be a family of topological spaces satisfying the \emph{gluing axiom}:
\begin{align}
X_{\lambda} \cap X_{\mu} \neq \emptyset \implies
\begin{aligned}
&\text{the relative topologies from } X_{\lambda} \text{ and } X_{\mu} \text{ coincide on } X_{\lambda} \cap X_{\mu},\\
&\text{and } X_{\lambda} \cap X_{\mu} = \mathcal{O}_{\lambda} \cap \mathcal{O}_{\mu}.
\end{aligned}
\label{eq:gluing_axiom}
\end{align}
Set $X := \bigsqcup_{\lambda \in \Lambda} X_{\lambda}$.
\begin{enumerate}[label=(\arabic*)]
\item Show that $\mathcal{O} := \{U \subset X : X_{\lambda} \cap U \in U_{\lambda} \forall \lambda \in \Lambda\}$ is a topology
on $X$.
\item Show that the relative topology on $X_{\lambda}$ from $(X, \mathcal{O})$ coincides with $\mathcal{O}_{\lambda}$ for each $\lambda \in \Lambda$.
\label{prob:2.2}
\end{enumerate}
\label{prob:2}
\end{problem} 

\begin{remark}
The gluing axiom \ref{eq:gluing_axiom} is necessary for \ref{prob:2.2} and to ensure $X_{\lambda} \in \mathcal{O}$. Such a topology $\mathcal{O}$ is uniquely determined
by the base $\bigcup_{\lambda \in \Lambda} \mathcal{O}_{\lambda}$.
\end{remark}

Let $(U, \phi)$ be any chart on $M$ and set
\[TU := \bigsqcup_{p \in U} T_pM = \{(p,v) : p \in U, v \in T_pM\} \subset TM.\]
Define $\tilde{\phi}: TU \to \phi(U) \times \mathbb{R}^n \subset \mathbb{R}^{2n}$ as follows:
For $p \in U$, $v \in T_pM$, let $\phi(p) = (x_1, \ldots, x_n)$ and $v = \sum_{i=1}^n v_i\partial_i$, then
\[\tilde{\phi}(p,v) := (x_1, \ldots, x_n, v_1, \ldots, v_n).\]

\begin{theorem}
\begin{enumerate}[label=(\arabic*)]
\item The family of topological spaces $\{TU\}$ satisfies the gluing axiom \ref{eq:gluing_axiom}.
Thus, $TM = \bigsqcup_{(U,\phi)} TU$ is a topological space with the topology definited in problem \ref{prob:2}.
\item $TM$ is a $2n$-dimensional manifold with atlas $\{(TU, \tilde{\phi})\}$.
\item The projection $\pi:TM \to M$ defined by $\pi(p,v) := p$ for $v \in T_pM \quad (p \in M)$ is a $C^{\infty}$ map.
\end{enumerate}
\end{theorem}